% titel: Kugel im Honig
% skript: Kap. Hydrodynamik

% Eine Stahlkugel sinkt in einem zähen Medium. Gefragt sind die wirkenden
% Kräfte (Stoke'sche Reibung, Auftrieb) und die Zeit bis fast zur
% Endgeschwindigkeit --- die klassische, durch lineare Reibung gedämpfte
% Fallbewegung. Dieser Teil erklärt die Physik von Grund auf; die Aufgabe
% wird NICHT durchgerechnet (das macht Teil II).

Eine Stahlkugel (Radius $r$, Masse $m$) sinkt in Honig (Viskosität $\eta$,
Dichte $\rho_f$). Gesucht sind die wirkenden Kräfte und die Zeit, bis die
Kugel einen Anteil $p$ (hier $p=99\%$) ihrer Endgeschwindigkeit erreicht hat.

\begin{werkzeugkasten}[Das physikalische Bild --- was passiert wirklich?]
Sobald die Kugel losgelassen wird, zieht die Schwerkraft nach unten. Aber sie
fällt \emph{nicht} frei und immer schneller wie im Vakuum, weil das zähe Honig
sie permanent bremst. Am sinkenden Körper greifen drei Kräfte an:

\begin{itemize}
  \item \textbf{Gewicht} $F_g = m g$ --- wirkt immer nach unten.
  \item \textbf{Auftrieb} $F_A=\rho_f V g$ --- nach Archimedes drückt das
        verdrängte Medium nach oben. Die Kugel verdrängt das Volumen
        $V=\frac43\pi r^3$; je dichter das Medium ($\rho_f$) und je grösser die
        Kugel ($V$), desto stärker der Auftrieb.
  \item \textbf{Strömungswiderstand (Stokes)} $F_R=6\pi\eta r\, v$ --- das Honig
        klebt an der Kugeloberfläche und zieht zurück gegen die Bewegung.
        \emph{Das Entscheidende:} dieser Widerstand wächst \textbf{linear mit
        der Geschwindigkeit} $v$. Langsam --- wenig Widerstand; schnell --- viel.
\end{itemize}
Vereinbart man \emph{nach unten als positiv}, so ist die Gesamtkraft
\[
  F_{\text{ges}} = m g - \rho_f V g - 6\pi\eta r\, v .
\]

\textbf{Warum es eine Endgeschwindigkeit gibt:} Am Anfang ist $v$ klein, die
Reibung vernachlässigbar, die Netto-Kraft zeigt nach unten --- die Kugel wird
schneller. Mit wachsender Geschwindigkeit wächst $6\pi\eta r\,v$, bis die
Reibung Gewicht und Auftrieb exakt aufhebt. Dann ist die Netto-Kraft null,
die Beschleunigung null, und die Kugel sinkt mit \textbf{konstanter
Geschwindigkeit} weiter: der Endgeschwindigkeit $v_\infty$.
\end{werkzeugkasten}

\begin{werkzeugkasten}[Bewegungsgleichung und die Notation der Übung]
2.~Newton: \emph{Masse $\times$ Beschleunigung = Summe der Kräfte}. Die
Beschleunigung ist die Ableitung der Geschwindigkeit, $\dot v=\mathrm dv/\mathrm dt$
(nach unten positiv):
\begin{equation}
  m\,\dot v = m g - \rho_f V g - 6\pi\eta r\, v .
\end{equation}
Teilt man durch $m$ und zieht zusammen, erhält man
\begin{equation}
  \dot v
  =\underbrace{g\Big(1-\frac{\rho_f V}{m}\Big)}_{=A}
   -\underbrace{\frac{6\pi\eta r}{m}}_{=B}\,v
  \;=\; A - B\, v .
\end{equation}
Dabei sind $A$ und $B$ \textbf{Konstanten}: $A$ ist der ``effektive Antrieb pro
Masse'' (Gewicht minus Auftrieb), $B$ die \textbf{Dämpfungsrate}. Die Grösse
$B$ hat die Einheit $\mathrm{s^{-1}}$, ist also ein Mass dafür, \emph{wie schnell
sich das System an die Endgeschwindigkeit angleich}.

\textbf{Warum die Übung $x'=A-Bx$ schreibt:} Das ist \emph{dieselbe} Gleichung,
nur ein anderer Buchstabe für dieselbe Grösse. Die Notation $x$ ist ein
allgemeiner Platzhalter für \emph{``die Grösse, die sich ändert''}. Bei der Kugel
ist diese Grösse die \textbf{Geschwindigkeit}, also ist $x=v$. Es gilt
\[
  x'=A-Bx \quad\Longleftrightarrow\quad v'=A-Bv ,
\]
gelesen als ``Änderungsrate = Antrieb minus Dämpfung $\times$ aktueller Wert''.
Weder am Rechenweg noch an der Lösung ändert sich etwas, wenn man $x$ durch
$v$ ersetzt.
\end{werkzeugkasten}

\begin{werkzeugkasten}[Lösung der DGL --- was sie bedeutet]
Die Gleichung $\dot v + Bv = A$ ist linear erster Ordnung mit konstanten
Koeffizienten. Mit dem integrierenden Faktor $e^{Bt}$ und der
Anfangsbedingung $v(0)=v_0$ erhält man
\begin{equation}
  v(t)=\frac{A}{B}+\Big(v_0-\frac{A}{B}\Big)e^{-Bt}.
\end{equation}
\begin{itemize}
  \item \textbf{Endwert (Endgeschwindigkeit):} Für $t\to\infty$ verschwindet
        $e^{-Bt}$, und es bleibt
        \[ v_\infty=\frac{A}{B}=\frac{\text{Antrieb}}{\text{Dämpfung}}. \]
        Genau dort gleichen Antrieb $A$ und Dämpfung $B v$ einander aus.
  \item \textbf{Einschwingen:} Von $v_0$ aus nähert sich die Geschwindigkeit
        \emph{exponentiell} $v_\infty$. Das Tempo bestimmt die
        \textbf{Zeitkonstante}
        \[ \tau=\frac{1}{B}=\frac{m}{6\pi\eta r}. \]
        Nach der Zeit $\tau$ ist die Restdifferenz $v_\infty-v(t)$ auf
        $e^{-1}\approx 37\%$ ihres Anfangswerts gesunken.
  \item \textbf{Zeit bis zum Bruchteil $p$:} Gesucht ist $t_*$ mit
        $v(t_*)=p\,v_\infty$. Mit $v_0=0$:
        \[ 1-e^{-Bt_*}=p \;\Rightarrow\;
           e^{-Bt_*}=1-p \;\Rightarrow\;
           \boxed{\;t_*=\frac1B\,\ln\frac{1}{1-p}\;} .
        \]
        Für $p=0.99$ wird $\ln\frac{1}{1-p}=\ln100$. Wichtig: $t_*$ hängt nur
        von $B$ (und $p$) ab, \emph{nicht} von $v_\infty$.
\end{itemize}
\end{werkzeugkasten}

\begin{werkzeugkasten}[Einheiten, Grössenordnung, Intuition]
\textbf{Einheiten:}
\[
 [B]=\frac{\mathrm{Pa\cdot m}}{\mathrm{kg}}
 =\frac{\mathrm{kg/(m\,s)\cdot m}}{\mathrm{kg}}=\mathrm{s^{-1}},
 \qquad (\text{also }) t_* \text{ in } \mathrm{s}.
\]
Umbauten: $1\,\mathrm{mPa\,s}=0.001\,\mathrm{Pa\,s}$;
$1\,\mathrm{kg/L}=1000\,\mathrm{kg/m^3}$.

\textbf{Intuition (Grenzfälle):}
\begin{itemize}
  \item \emph{Schwerere, dichtere Kugel} (grösseres $A$): grössere $v_\infty$.
        Für $\rho_f V = m$ (Kugel hat dieselbe Dichte wie das Medium) ist
        $A=0$ und $v_\infty=0$ --- sie ``schwimmt''.
  \item \emph{Zäheres Medium oder grössere Kugel} (grösseres $B$): kleinere
        $v_\infty$, aber auch kürzeres $\tau$ --- die (kleinere)
        Endgeschwindigkeit wird in kürzerer Zeit erreicht.
  \item Zusammengefasst: $v_\infty\propto (m-\rho_f V)/\eta$ und
        $t_*\propto m/(\eta r)$.
\end{itemize}
Stokes gilt für kleine Reynoldszahlen (schleichende Strömung); bei hoher
Viskosität (Honig) ist das erfüllt.
\end{werkzeugkasten}

\begin{werkzeugkasten}[Vorgehen in allgemeiner Form]
\begin{enumerate}
  \item Die drei Kräfte benennen und die Richtung festlegen (nach unten $+$):
        $F_g - F_A - F_R$.
  \item $F_A=\rho_f\frac43\pi r^3 g$ und $F_R=6\pi\eta r\,v$ einsetzen und auf
        die Form $\dot v = A - B v$ bringen; $A$ und $B$ herauslesen.
  \item Endgeschwindigkeit $v_\infty=A/B$ bilden.
  \item Mit $v(0)=0$ die Lösung $v(t)=v_\infty(1-e^{-Bt})$ nutzen und aus
        $1-e^{-Bt_*}=p$ die Zeit $t_*=\ln(1/(1-p))/B$ bestimmen.
  \item Einheiten jedes Zwischenergebnisses prüfen.
\end{enumerate}
\end{werkzeugkasten}

\begin{achtung}[Typische Fallen]
\begin{itemize}
  \item Vorzeichen: Auftrieb und Reibung wirken \emph{entgegen} der Bewegung
       (nach oben), Gewicht nach unten.
  \item Dichte: den Auftrieb mit der \emph{Honig}-Dichte $\rho_f$ rechnen
       (nicht mit der der Kugel).
  \item Einheiten: $100\,\mathrm{mPa\,s}=0.1\,\mathrm{Pa\,s}$;
       $1\,\mathrm{kg/L}=10^3\,\mathrm{kg/m^3}$ --- Faktor 1000 nicht vergessen.
  \item Anfangswert: für $v(0)=0$ ist die Konstante in der Lösung
       $C=-A/B$; fehlt sie, stimmt $v(t)$ nicht.
  \item $v_\infty$ in $\mathrm{m/s}$ erst am Schluss in $\mathrm{mm/s}$
       umrechnen.
\end{itemize}
\end{achtung}
